\begin{abstract}
  An increasing number of companies build their business on
  distributed Web applications. Hosting providers respond to that
  demand and made it easier to deploy systems that spread across
  multiple services. However, this trend has outpaced the development
  of adequate debugging tools and developers still have to rely on an
  improvised patchwork of symbolic debuggers and printf debugging to
  find failure causes.

  Consequently, observing and debugging failures in distributed
  applications is laborious and time consuming. As multiple processes
  execute in parallel, developers cannot systematically stop and
  inspect the entire system. Although processes are implemented in
  different languages, and often tightly coupled, debugging tools that
  span across processes and languages are practically
  non-existent. Instead, developers have to use different debugging
  techniques to inspect each process sequentially. This makes it hard
  to reason about causality, because the network communication changes
  when debugging a process. That may lead to failures that occur in
  deployed systems, but that developers cannot reproduce locally.

  This work presents an integrated method to support debugging of
  distributed systems in four steps: First, we continually trace
  network communication to have accurate data about event order if a
  failure occurs in a deployed system. Second, we compare
  communication patterns heuristically to identify anomalous
  differences and guide developers in search of failure causes. Third,
  we replay failing network schedules to reproduce failures in a local
  test environment. Finally, we refine the traced data during replays
  and connect network events to code locations. We present network
  events and code within one tool, so developers can use both levels
  of abstraction for debugging.

  We show how our approach can be used to inspect failures in
  distributed systems that would be hard to debug with traditional
  approaches.
\end{abstract}

\begin{zusammenfassung}
  Die Zahl der Firmen, die als Teil ihres Geschäftes verteilte Web
  Anwendungen entwickeln, nimmt zu. Anbieter zur Bereitstellung von
  Servern bedienen die steigende Nachfrage, und erleichtern den
  Einsatz von Systemen die sich über mehrere Dienste hinweg
  erstrecken. Dieser Trend jedoch schreitet schneller voran, als die
  Diagnoseprogramme für diesen Einsatzbereich weiterentwickelt werden,
  weshalb Entwickler noch immer mit symbolischen Fehlersuchwerkzeugen
  und Protokollierung arbeiten müssen, um Fehler zu finden.

  Folglich ist das Beobachten und Beheben von Fehlern in verteilten
  Systemen noch immer mühselig und zeitaufwendig. Da mehrere Prozesse
  parallel laufen, können Entwickler nicht das gesamte System zur
  Untersuchung anhalten. Obwohl Prozesse in mehreren Sprachen
  implementiert, und eng gekoppelt sind, gibt es praktisch keine
  Werkzeuge, die über Prozess- und Sprachgrenzen hinweg
  funktionieren. Stattdessen müssen Entwickler verschiedene
  Fehlersuchtechniken für jeden Prozess einzeln anwenden. Das
  erschwert es ihnen, kausale Zusammenhänge zu verstehen, da sich das
  Kommunikationsverhalten ändert, wenn einzelne Prozesse untersucht
  werden. Das kann dazu führen, dass Fehler, die im Produktivsystem
  auftreten, von Entwicklern lokal nicht reproduziert werden können.

  In dieser Arbeit wird eine Methode zur Fehlersuche in verteilten
  Systemen präsentiert, welche die folgenden vier Schritte umfasst:
  Zunächst wird fortwährend die Kommunikation im Netzwerk
  aufgezeichnet, um im Falle eines Fehlers im Produktivsystem den
  Entwicklern akkurate Ereignisdaten bereitzustellen. Daraufhin werden
  diese Daten mithilfe einer Heuristik auf Anomalien in den
  Kommunikationsmustern hin untersucht. Weiterhin werden diese Daten
  zur Wiedergabe der Kommunikation in einem Entwicklungssystem
  genutzt. Schließlich werden die aufgezeichneten Daten während der
  Wiedergaben verfeinert, und so die Netzwerkdaten mit dem Quelltext
  assoziiert. Beides wird in einem Werkzeug präsentiert, sodass
  Entwickler beide Abstraktionsebenen zur Fehlersuche nutzen können.

  Wir zeigen wie unsere Methode genutzt werden kann, um Fehler in
  verteilten Systemen zu finden, die mit traditionellen Ansätzen nur
  schwierig untersucht werden können.
\end{zusammenfassung}
 

%%% Local Variables: 
%%% mode: latex
%%% TeX-master: "thesis"
%%% End: 
